% ============================================================
% LLM-BLM 実験仮説（設計書）
% コンパイル：tectonic -X compile LLM-BLM_実験仮説.tex   （XeTeX 系。このフォルダで実行する）
% 実験の前に書いた設計書。結果は『LLM-BLM_実験結果レポート』にまとめる。
% ============================================================
\documentclass[xelatex,ja=standard,jafont=haranoaji,a4paper,10pt,openany]{bxjsreport}
\setpagelayout*{left=24mm,right=24mm,top=24mm,bottom=26mm,footskip=11mm}

\usepackage{amsmath,amssymb,mathtools,bm}
\usepackage{booktabs,array,multirow,longtable}
\usepackage{graphicx}
\usepackage{xcolor}
\usepackage{tikz}
\usetikzlibrary{arrows.meta,positioning,calc,fit}
\usepackage{listings}
\usepackage{enumitem}
\usepackage{caption}
\usepackage[most]{tcolorbox}
\usepackage[hidelinks,bookmarksnumbered]{hyperref}
\pdfstringdefDisableCommands{\def\Omega{Ω}\def\tau{τ}\def\Sigma{Σ}\def\Pi{Π}\def\mu{μ}\def\delta{δ}\def\T{ᵀ}}

% ---- 色（総合報告書・BLM解説・ビューとリスクのコツと同じ）
\definecolor{c1}{HTML}{2A78D6}
\definecolor{c2}{HTML}{EB6834}
\definecolor{c3}{HTML}{1BAF7A}
\definecolor{c4}{HTML}{EDA100}
\definecolor{c7}{HTML}{6250D6}
\definecolor{ink2}{HTML}{52514E}
\definecolor{soft1}{HTML}{E4EEF9}
\definecolor{soft2}{HTML}{FBEEDF}
\definecolor{soft3}{HTML}{E3F3E9}
\definecolor{soft4}{HTML}{EEE8FA}
\definecolor{soft5}{HTML}{FDF3DA}

% ---- よく使う表記
\newcommand{\code}[1]{\texttt{#1}}
\newcommand{\SP}{S\&P500}
\newcommand{\T}{^{\top}}
\newcommand{\muBL}{\mu_{\mathrm{BL}}}
\newcommand{\wmkt}{w_{\mathrm{mkt}}}
\newcommand{\one}{\mathbf{1}}
\newcommand{\IC}{\mathrm{IC}}
\newcommand{\IR}{\mathrm{IR}}
\newcommand{\TE}{\mathrm{TE}}
\DeclareMathOperator{\diag}{diag}
\newcommand{\figref}[1]{図~\ref{#1}}
\newcommand{\tabref}[1]{表~\ref{#1}}
\newcommand{\chapref}[1]{第~\ref{#1}~章}
\newcommand{\secref}[1]{\ref{#1}~節}

\setlist[itemize]{leftmargin=1.6em,itemsep=0.2em,topsep=0.4em}
\setlist[enumerate]{leftmargin=1.8em,itemsep=0.2em,topsep=0.4em}
\captionsetup{font=small,labelfont=bf,skip=6pt}
\renewcommand{\arraystretch}{1.2}
\setlength{\LTpre}{6pt}
\setlength{\LTpost}{6pt}

% ---- 章の見出しを小さくする（ビューとリスクのコツと同じ）
\makeatletter
\def\@makechapterhead#1{\vspace*{0.3\Cvs}{\parindent\z@\raggedright\normalfont
  \ifnum\c@secnumdepth>\m@ne\if@mainmatter{\Large\headfont\@chapapp\thechapter\@chappos}\par\nobreak\vskip0.4\Cvs\fi\fi
  \interlinepenalty\@M\LARGE\headfont #1\par\nobreak\vskip 1.5\Cvs}}
\def\@makeschapterhead#1{\vspace*{0.3\Cvs}{\parindent\z@\raggedright\normalfont
  \interlinepenalty\@M\LARGE\headfont #1\par\nobreak\vskip 1.5\Cvs}}
\makeatother

% ---- コード
\renewcommand{\lstlistingname}{リスト}
\lstset{basicstyle=\ttfamily\footnotesize,breaklines=true,columns=fullflexible,keepspaces=true,
  showstringspaces=false,frame=single,framerule=0.3pt,rulecolor=\color{black!25},backgroundcolor=\color{black!3},
  xleftmargin=3pt,xrightmargin=3pt,aboveskip=8pt,belowskip=8pt,captionpos=t}

% ---- 囲み。事実・仮説・予想を色で分ける
\newtcolorbox{keybox}[1][要点]{enhanced,breakable,colback=soft1,colframe=c1,boxrule=0.5pt,arc=2pt,
  left=8pt,right=8pt,top=6pt,bottom=6pt,title={#1},fonttitle=\sffamily\bfseries\small,coltitle=white}
\newtcolorbox{notebox}[1][注意]{enhanced,breakable,colback=soft2,colframe=c2,boxrule=0.5pt,arc=2pt,
  left=8pt,right=8pt,top=6pt,bottom=6pt,title={#1},fonttitle=\sffamily\bfseries\small,coltitle=white}
\newtcolorbox{factbox}[1][資料で確認した事実]{enhanced,breakable,colback=soft3,colframe=c3,boxrule=0.5pt,arc=2pt,
  left=8pt,right=8pt,top=6pt,bottom=6pt,title={#1},fonttitle=\sffamily\bfseries\small,coltitle=white}
\newtcolorbox{hypbox}[1][仮説]{enhanced,breakable,colback=soft4,colframe=c7,boxrule=0.5pt,arc=2pt,
  left=8pt,right=8pt,top=6pt,bottom=6pt,title={#1},fonttitle=\sffamily\bfseries\small,coltitle=white}
\newtcolorbox{predbox}[1][実験前の予想]{enhanced,breakable,colback=soft5,colframe=c4!85!black,boxrule=0.5pt,arc=2pt,
  left=8pt,right=8pt,top=6pt,bottom=6pt,title={#1},fonttitle=\sffamily\bfseries\small,coltitle=white}

% ---- 図の部品
\tikzset{
  box/.style={draw=#1,fill=#1!8,rounded corners=2pt,align=center,font=\footnotesize,inner sep=4pt,
              minimum height=10mm,text width=23mm},
  arr/.style={-{Stealth[length=2mm]},color=ink2,line width=0.6pt},
  lab/.style={font=\scriptsize,color=ink2},
}

% ---- 仮説ごとの 13 項目の表
\newenvironment{plan}{\small\begin{longtable}{p{3.3cm}p{11.6cm}}\toprule\textbf{項目} & \textbf{内容} \\\midrule\endfirsthead
\toprule\textbf{項目} & \textbf{内容} \\\midrule\endhead\bottomrule\endfoot}{\end{longtable}}

\title{LLM-BLM 実験仮説\\[6pt]
  \large 既存の資料と実装の調査にもとづく改善仮説と、その検証の計画}
\author{実験の設計書（実験の前に作成）}
\date{2026 年 10 月 16 日}

\begin{document}
\maketitle
\setcounter{tocdepth}{1}
{\small\tableofcontents}

% ============================================================
\chapter*{この資料について}
\addcontentsline{toc}{chapter}{この資料について}

本書は、LLM-BLM プロジェクトについて、既存の資料・コード・実験結果を調べたうえで立てた改善仮説と、その検証の計画をまとめたものである。\textbf{実験を始める前に書き、Git に記録してから実験に入る。}評価の指標と、仮説を支持する・しないの条件は、結果を見る前にここで決めておく。結果を見てから都合のよい指標を選ぶことを防ぐためである。

研究の目的は、LLM を使うこと自体を正当化することではない。\textbf{LLM が、既存の Black-Litterman モデル（BL）や単純な統計の方法に対して、本当に追加の価値を生むか}を確かめることである。

\section*{書き分け}

本書では、次の 3 つを囲みの色で区別する。

\begin{factbox}[資料で確認した事実（緑）]
既存の資料・コード・データで、この設計書を書く前に確かめたこと。出典（資料の章や、コードのファイル）を添える。
\end{factbox}
\begin{hypbox}[新しく立てる仮説（紫）]
本書で新しく提案し、実験で確かめる主張。
\end{hypbox}
\begin{predbox}[実験前の予想（黄）]
結果について、実験の前に持っている見込み。外れても、そのまま記録する。
\end{predbox}

\section*{設計のレビュー}

最初の版を、2 つの LLM（OpenAI の gpt-5.5、Google の gemini-3.1-pro-preview）に読ませ、問題点を挙げさせた。指摘を自分で確かめ、次の点を直した。H02〜H04 が既存のデータの再分析であることの明記、判断日の翌取引日に売買する感度分析、H07 の銘柄の選び方（後の在籍を条件にしない）、「有意でない＝効果なし」としない判定（意味のある最小の効果を IC 0.03 に決める）、$b_t=0$ の扱い、順位ビューを同じ BL の式で比べる形への変更、H02 の主な検定を残差の IC に変更（多重共線性）、置換を 200 回に増やす、決算発表文の HTML を取り除く、H05 のなぞりの判定を検定に変更。指摘のうち「仮説ごとに補正せず単独で評価すべき」は採らず、補正の前と後の両方を書くことにした。

\section*{調べた資料}

\begin{center}
\small
\begin{tabular}{p{5.4cm}p{8.8cm}}
\toprule
資料 & 読んだ範囲 \\
\midrule
『LLM-BLM 総合報告書』（\code{報告書/LLM-BLM\_総合報告書.pdf}、原稿は \code{資料いろいろ/LLM-BLM\_総合報告書.tex}） & 本文と付録（プロンプトの全文、月別の数値）を通読した。依頼文の「総合報告書(2).pdf」に当たるファイルは見つからず、この版を使った \\
『BLM ビューとリスクのコツ』（\code{LLM-BLMビューとリスクのコツ/}） & $\tau$、$\Omega$、当たり具合、積み上がり、ビューの形、TE、本プロジェクトへの当てはめ、質問の設計、コツの一覧の章を通読した。ほかの章は要点の囲みと表を読んだ。計算スクリプトの E、E2 節（$\Omega$ の比較）を読んだ \\
『BLM 解説』（\code{勉強用/}） & 構成と式の記法を確認した \\
実験のスクリプト（\code{LLM-BLM\_実験/}） & 実験 D〜G のスクリプトと結果のファイルを確認した \\
コード（\code{blm/}、\code{llm/}、\code{marketdb/}、\code{run\_*.py}） & すべて読んだ \\
データベース（\code{data/blm\_2015.sqlite}） & テーブル、LLM の実行、比率の計算の記録を確認した \\
\bottomrule
\end{tabular}
\end{center}

% ============================================================
\chapter{既存の研究で分かっていること}\label{ch:known}

\section{総合報告書の主な結果}

\begin{factbox}[総合報告書（2026 年 10 月 13 日版）で報告されていること]
\begin{itemize}
\item 元論文の計算では、BL が均衡から LLM の予想へ動かす割合（採用率）は中央値で 99.87\%（最小 99.75\%）。均衡リターンをすべて 0 にしても比率は平均 92\% 重なる（第 III 部、問題 1）。
\item LLM の予想と、渡した月の平均リターンの相関は平均 0.89（最小 0.80、最大 0.95）。回帰の傾きは 0.77（第 6 章）。
\item 予想と翌月の実際のリターンの順位相関は 30 か月の平均で $+0.05$。プラスの月は 14 か月（第 6 章、付録の月別の表）。
\item 回答の分散 ÷ 実際の予測誤差の目安は中央値で 0.05〜0.15（問題 3）。
\item 元の最適化（SciPy の SLSQP、標準の設定）は、改善できる幅の 95\% で止まり、実質 9.9 銘柄。同じ式を二次計画の解法で最後まで解くと 4.2 銘柄（問題 4）。
\item 均衡リターンの計算で、日次のリターンから年率の金利 0.02 をそのまま引いている。ただし採用率が 99.9\% なので結果にはほぼ効かない（問題 5）。
\item 共分散は判断月の約 20 日分の標本共分散で、約 488 銘柄に対して階数が足りない（問題 6）。
\item 直した計算（逆最適化の均衡、252 日の Ledoit--Wolf、Clarabel で最後まで解く）では、比率が実質 1.7 銘柄に集中する（第 7 章）。
\item 入力を要約した指標に変える（実験 D）、比率を明言する（E）、3 ターンにする・5 段階で答えさせる（F）、銘柄の数字を渡さずリスクの種類ごとに聞く（G）のどれでも、答えは渡した数字の中の直近の値動きに沿った。
\item 結果へのフィードバックを行う LLM は「未着手」である（第 1 章の現在の状態の表）。\textbf{コードにも実装はない。}
\end{itemize}
\end{factbox}

\section{BLM ビューとリスクのコツの主な結果}

\begin{factbox}[ビューとリスクのコツ（2026 年 10 月）で報告されていること]
\begin{itemize}
\item $\Omega$ を $\tau$ に比例させると、事後の期待リターンから $\tau$ は消える。$\Omega$ を外から決めると、$\tau$ は採用率を直接決める係数になる。$\tau$ でトラッキングエラー（TE）を合わせると、「どれだけずれてよいか」と「ビューをどれだけ信じるか」を混ぜてしまう。$\Omega$ は当たり具合で、TE は最適化の制約で決めるのがよい（第 4、18 章）。
\item 市場版の BL で $\omega=\sigma^2(1/\IC^2-1)$ と置くと、Grinold の式 $\alpha=\IC\cdot\sigma\cdot z$ と一致し、採用率は $\IC^2$ になる（第 8 章）。
\item ビューが多いのに $\Omega$ を対角にすると、外れの共通部分が無視され、確信が積み上がる。本プロジェクトの LLM の外れの共通相関は約 0.21、全体の水準の予想と実際の相関は $-0.43$（第 8b 章）。
\item LLM の当たり具合：順位相関の平均 $+0.052$（$t$ 値 1.51）、月ごとに平均を引いてまとめた傾き $+0.074$、回答の分散は実際の外れの 2 乗の中央値の 0.124 倍（第 19 章）。
\item $\Omega$ を 7 通りに変えた 30 か月の計算（第 19 章の表）：回答の分散では採用率 91.3\%・実質 1.7 銘柄・TE 50.9\%。傾きで較正し $\Sigma$ に比例させると採用率 5.6\%・実質 17.6 銘柄・TE 26.0\%。さらに事前の TE を 3\% 以内にすると実質 41.3 銘柄・実績の TE 5.3\%・IR $-0.02$。\textbf{どの方法も、超過リターンの $t$ 値の絶対値は 0.6 未満}だった。
\item 2015 年からの銘柄の指標の IC（\SP{} の構成銘柄、129〜139 か月）は、どれも絶対値 0.03 未満で、$t$ 値の絶対値は 1.2 未満（第 13 章）。
\item 30 か月では、IC 0.05 程度の予想の価値は確かめられない。$t$ 値 2 で確かめるには約 55 か月が要る（第 8 章）。
\end{itemize}
\end{factbox}

\section{まだ確かめられていないこと}

資料が「次に行うこと」に挙げ、まだ行われていないものを並べる。

\begin{center}
\small
\begin{tabular}{p{6.0cm}p{8.2cm}}
\toprule
未検証の課題 & 出典 \\
\midrule
文章だけを渡したとき、LLM の答えは月によって変わるか。数字に沿っていないか & 総合報告書 第 V 部の総合考察、次に行うこと 1 \\
統計モデル（回帰）を比べる相手として作り、LLM が何を足しているかを示す & 総合報告書 第 10 章、次に行うこと 4 \\
銘柄ごとの文章の情報源（決算の説明の文章など）を公表日つきで手に入れる & 総合報告書 次に行うこと 5 \\
結果へのフィードバックを行う LLM & 総合報告書 第 1 章、次に行うこと 7 \\
較正の傾き $b$ を、推定の誤差に応じて 0 に寄せて使う & ビューとリスクのコツ 第 8 章（「簡単のため縮めずに使った」） \\
売買のコストを入れた比較 & 両資料の限界の節 \\
較正の仕組みを \code{blm/} に組み込む & ビューとリスクのコツ 第 19 章の提案 \\
\bottomrule
\end{tabular}
\end{center}

% ============================================================
\chapter{現行のアーキテクチャ}\label{ch:arch}

\section{データの流れ}

\begin{figure}[htbp]
\centering
\begin{tikzpicture}[node distance=5mm and 4mm]
\node[box=c3] (src) {外部データ\\{\scriptsize Wikipedia、yfinance、EDGAR（数値）}};
\node[box=c1,right=of src] (mdb) {\code{marketdb/}\\{\scriptsize 取得・加工・指標}};
\node[box=c4,right=of mdb,minimum height=16mm] (db) {SQLite\\{\scriptsize \code{blm\_2015.sqlite}}};
\node[box=c1,above right=2mm and 6mm of db] (llm) {\code{run\_llm.py}・\code{llm/}\\{\scriptsize 当月の日次リターン $\to$ 30 個の予想}};
\node[box=c2,right=of llm] (api) {OpenAI API\\{\scriptsize gpt-4o-mini, $T=1$, $n=30$}};
\node[box=c1,below right=2mm and 6mm of db] (blm) {\code{run\_blm.py}・\code{blm/}\\{\scriptsize $\Pi$、$\Sigma$、$\muBL$、最適化}};
\node[box=c4,right=of blm] (ev) {評価\\{\scriptsize ビュー \code{portfolio\_returns}}};
\draw[arr] (src) -- (mdb);
\draw[arr] (mdb) -- (db);
\draw[arr,{Stealth[length=2mm]}-{Stealth[length=2mm]}] (db) |- (llm);
\draw[arr] (llm) -- (api);
\draw[arr,{Stealth[length=2mm]}-{Stealth[length=2mm]}] (db) |- (blm);
\draw[arr] (blm) -- (ev);
\end{tikzpicture}
\caption{現行のデータの流れ。\code{llm/} と \code{blm/} はデータベースを通してだけつながる。結果へのフィードバックを行う LLM はない。}
\label{fig:arch}
\end{figure}

\section{処理ごとの確認}

\begin{longtable}{p{3.4cm}p{11.5cm}}
\toprule
\textbf{処理} & \textbf{現行の実装（コードで確認）} \\
\midrule
\endfirsthead
\toprule
\textbf{処理} & \textbf{現行の実装（コードで確認）} \\
\midrule
\endhead
\bottomrule
\endfoot
1. データの取得元 & 構成銘柄と業種は Wikipedia、株価と株式数と指数は yfinance、財務諸表は EDGAR の XBRL（\code{marketdb/sources/}）。期間は 2014-01-02〜2026-09-30（株価）。銘柄の指標は 2015-02〜2026-09 の 140 か月（\code{stock\_metrics}） \\
2. LLM への入力 & ティッカー、社名、セクター、サブ業種、判断月の日次リターンの列（約 20 個、小数）。日付は渡さない（\code{llm/prompts.py}、\code{llm/context.py}） \\
3. LLM の出力 & 翌月の平均日次リターン（小数）を 1 銘柄 30 個。構造化出力で \code{\{"expected\_return": 数値\}}。文章や信頼度は出させない（\code{llm/client.py}）。実行 1 は gpt-4o-mini-2024-07-18、温度 1.0、判断日 2024-03-28〜2026-08-31 の 30 か月、のべ 14,688 件 \\
4. $q$、$P$、$\Omega$ への変換 & $q$ = 30 個の平均 $-$ 日次の無リスク金利、$P=I$（全銘柄の絶対ビュー）、$\Omega=\diag(\text{30 個の母分散})$（\code{blm/strategies.py}、ビュー \code{llm\_view\_stats}） \\
5. $\Pi$、$\Sigma$、$\tau$ & $\Sigma$ は過去 252 取引日の日次リターンの Ledoit--Wolf 推定（\code{blm/covariance.py}）。$\Pi=\delta\Sigma\wmkt$、$\delta=2.5$（\code{blm/equilibrium.py}）。$\tau=0.025$。$\wmkt$ は判断日の時価総額（終値 $\times$ 発行済み株式数、浮動株の調整なし） \\
6. 事後と最適化 & $\muBL=\Pi+\tau\Sigma(\tau\Sigma+\Omega)^{-1}(q-\Pi)$（\code{blm/black\_litterman.py}）。$\max_w w\T\muBL-\frac{\delta}{2}w\T\Sigma w$、$\one\T w=1$、$w\ge0$ を CVXPY と Clarabel で解く。最適でなければ例外を出す（\code{blm/optimizer.py}）。TE の制約、売買のコストはない \\
7. バックテスト & 判断日の終値で買い、翌月の最後の取引日まで持つ。配当込みの調整後終値で計算（ビュー \code{portfolio\_returns}）。比べる相手は時価総額加重。指標は \code{blm/evaluation.py} \\
8. DB への保存 & LLM の実行条件（\code{llm\_runs}：モデル、システムプロンプト、温度、回答数）、全回答（\code{llm\_views}）、比率の実行条件（\code{portfolio\_runs}：係数を JSON で）と比率（\code{portfolio\_weights}）。予想と実現値の対応は、判断日と銘柄で株価に結合して作る。\textbf{ユーザープロンプトの型と、使った入力のまとまりは保存されていない} \\
9. フィードバックの LLM & \textbf{実装されていない。}総合報告書の計画（段階 2）では「数字で要因を分解し、その結果と当時のレポートを LLM に渡して説明させる」とされている \\
\end{longtable}

\section{仮説の検証のために変える場所}

\begin{center}
\small
\begin{tabular}{p{4.2cm}p{10.0cm}}
\toprule
変える要素 & 変える場所 \\
\midrule
ビューの形（$q$、$P$） & \code{blm/black\_litterman.py} に、一般の $P$ と行列の $\Omega$ を受け取る関数を足す \\
$\Omega$ の決め方と較正 & 較正の関数を \code{blm/calibration.py} に足す（判断日より前のデータだけを使う） \\
TE の制約、売買のコスト & \code{blm/optimizer.py} に TE つきの最適化を足す。コストは評価の側で引く \\
LLM への質問 & \code{experiments/} の下に実験ごとのスクリプトを置く。回答と条件は実験のフォルダに保存し、既存の DB には書き込まない \\
統計モデルの対照 & \code{stock\_metrics} を読む実験のスクリプト \\
\bottomrule
\end{tabular}
\end{center}

既存の \code{.env} の決まり（コードに既定値を置かない）を崩さないため、新しい関数は設定を読まない純粋な関数にし、係数は実験のスクリプトの中で定める。既存の関数と \code{run\_*.py} の動きは変えない。

% ============================================================
\chapter{数式と、整合性を確かめる項目}\label{ch:math}

\section{標準の式}

均衡の期待超過リターン（逆最適化）と、BL の事後の期待リターンは次のとおりである。
\begin{align}
\Pi&=\delta\,\Sigma\,\wmkt\label{eq:pi}\\
\muBL&=\left[(\tau\Sigma)^{-1}+P\T\Omega^{-1}P\right]^{-1}\left[(\tau\Sigma)^{-1}\Pi+P\T\Omega^{-1}q\right]\label{eq:bl1}\\
&=\Pi+\tau\Sigma P\T\left(P\tau\Sigma P\T+\Omega\right)^{-1}(q-P\Pi)\label{eq:bl2}
\end{align}
$N$ を銘柄数、$K$ をビューの数とすると、$\Sigma$ は $N\times N$、$P$ は $K\times N$、$q$ は $K$ 次元、$\Omega$ は $K\times K$ である。現行の実装は $P=I$（$K=N$）で、式~\eqref{eq:bl2} の形を \code{numpy.linalg.solve} で解いている。

\section{単位}

現行の「直した計算」では、$\Sigma$ は日次リターンの共分散、$\Pi$ は日次の超過リターン、$q$ は「翌月の平均日次リターン $-$ 日次の無リスク金利」、$\Omega$ は「30 個の回答の分散」で、平均日次リターンの分散である。$\Pi$ と $q$ と $\tau\Sigma$ はどれも「日次の平均リターン」の単位で、そろっている。

ただし、$\Omega$ が表すのは「回答のぶれ」であり、「平均日次リターンの予想の外れ」ではない。この点は計算の誤りではなく、研究の論点（仮説 H04）として扱う。

\section{確かめる項目}

\begin{enumerate}
\item \textbf{次元}：$P$ が一般の形（相対ビュー、$K<N$）でも式~\eqref{eq:bl1} と式~\eqref{eq:bl2} が一致すること。
\item \textbf{ビューがないとき}：$\Pi$ をそのまま最適化に入れると、比率が $\wmkt$ に戻ること（30 か月すべて）。
\item \textbf{$\Omega$ の作り方}：ビュー \code{llm\_view\_stats} の分散が、30 個の回答の母分散と一致すること。
\item \textbf{逆行列の安定性}：$\tau\Sigma+\Omega$ の条件数と、\code{solve} の残差。
\item \textbf{採用率}：銘柄ごとの $(\muBL-\Pi)/(q-\Pi)$ の中央値が、資料の値（元の計算 99.87\%、直した計算 91.3\%）を再現すること。
\item \textbf{最適化の収束}：Clarabel の状態が「最適」であること。目的関数の値を、別の解き方と比べること。元の SLSQP の停止の理由（成功の印、反復の回数）と、最適値からの差を記録すること。
\item \textbf{再現}：同じ DB と同じコードで、資料の数値（相関 0.89、順位相関 $+0.05$、実質銘柄数）を出し直せること。
\end{enumerate}

計算の誤りが見つかった場合は、研究の仮説と混ぜず、まず計算の問題として記録して直す。

% ============================================================
\chapter{仮説の全体像}\label{ch:overview}

\section{仮説の一覧}

\begin{center}
\small
\begin{tabular}{p{0.9cm}p{1.2cm}p{9.8cm}p{2.0cm}}
\toprule
ID & テーマ & 名称 & 費用の見込み \\
\midrule
H01 & D & 現行の計算の再現と整合性の確認（計算の検証。研究の仮説ではない） & 0 \\
H02 & F、A & LLM の数値予想は、入力の単純な変換と統計モデルを超える情報を持つか & 0 \\
H03 & B & $q$ と $P$：絶対ビュー、相対ビュー、順位ビュー & 0 \\
H04 & C、D & $\Omega$：回答の分散、較正、縮小、置換による対照 & 0 \\
H05 & A & 市場とセクター：文章だけ・数字だけで聞く（実験 G の拡張） & 約 0.2 ドル \\
H06 & A、F & 銘柄ごとの決算発表文から LLM が情報を取り出し、統計で数値化する & 約 5 ドル \\
H07 & E & フィードバックの LLM：何を渡すと何が変わるか & 約 3 ドル \\
\bottomrule
\end{tabular}
\end{center}

テーマの記号は、依頼文の A〜F（A：LLM に何を予測させるか、B：$q$ と $P$、C：$\Omega$、D：BLM と最適化、E：フィードバック、F：統計モデルとの比較）である。費用は gpt-4o-mini の単価（入力 100 万トークンあたり 0.15 ドル、出力 0.60 ドル）による見込みで、実績は結果のレポートに書く。総合報告書の実行 1（14,688 リクエスト、各 30 回答）が約 1.75 ドルだったことと、この単価は合っている。

\section{実験の順番と、その理由}

依頼文の暫定の順番（検証 → ビューの変換 → $q$・$P$・$\Omega$ → 較正 → フィードバック → 統計モデルとの比較）から、次のように変えた。

\begin{enumerate}
\item \textbf{H01（計算の検証）を最初に置く。}後の比較のすべてが、この計算に乗るためである。
\item \textbf{統計モデルとの比較（H02）を 2 番目に上げる。}依頼文では最後だったが、統計モデルの予想は H03 以降の「比べる相手」になる。先に作らないと、各仮説の対照群がそろわない。
\item \textbf{H03（$q$、$P$）と H04（$\Omega$）は同じ評価の仕組みを使う。}H03 では $\Omega$ を、H04 では $q$ を固定する。
\item \textbf{LLM を新しく呼ぶ実験（H05〜H07）は後に置く。}API を使う前に、評価の仕組みを H01〜H04 で確かめておくためである。H05 は小さく安いので先に行い、規模の大きい H06 の前の動作確認も兼ねる。
\end{enumerate}

\section{ブランチの作り方}

\begin{center}
\small
\begin{tabular}{ll}
\toprule
対象 & ブランチ \\
\midrule
設計書（本書） & \code{research/llm-blm-hypotheses} \\
H01 & \code{experiment/H01-verification} \\
H02 & \code{experiment/H02-statistical-baseline} \\
H03 & \code{experiment/H03-view-shape} \\
H04 & \code{experiment/H04-omega-calibration} \\
H05 & \code{experiment/H05-text-vs-numbers} \\
H06 & \code{experiment/H06-earnings-text} \\
H07 & \code{experiment/H07-feedback-llm} \\
結果の報告書 & \code{research/llm-blm-results-report} \\
\bottomrule
\end{tabular}
\end{center}

本書は \code{research/llm-blm-hypotheses}（\code{db-from-2015} から分岐）に置く。各仮説のブランチは、1 つ前の仮説のブランチから順に分岐する（H01 は本書のブランチから）。共通の評価の仕組みを H01 で作り、後の仮説で使い回すためである。最終の報告書は、最後の実験のブランチから \code{research/llm-blm-results-report} を分岐して作る。

% ============================================================
\chapter{共通の評価方法}\label{ch:eval}

\section{期間と時点}

\begin{itemize}
\item \textbf{評価の判断日}：2024-03-28〜2026-08-31 の 30 か月の、各月の最後の取引日（実行 1 と同じ）。予測の対象は、判断日の翌日から翌月の最後の取引日まで。
\item \textbf{学習に使う期間}：統計モデルは 2015-02〜（判断日の前に結果が確定した月まで）。ウォークフォワード（判断日ごとに、その時点で結果が分かっている月だけで学習・較正する）で行う。判断日 $d_i$ の時点では、判断日 $d_{i-1}$ の予想の結果（$d_i$ で終わる月のリターン）までが分かっている。
\item \textbf{未来の情報を使わない決まり}：入力、較正、銘柄の選択に、判断日より後の情報を使わない。決算発表文は、提出日が判断日より\textbf{前}（同じ日は除く）のものだけを使う（同じ日の取引終了後の提出がありうるため）。
\item \textbf{LLM のモデル}：銘柄や相場の予想に使うのは gpt-4o-mini（知識は 2023 年 10 月まで）だけにする。新しいモデル（gpt-5.x、Gemini 3.x）は評価期間の相場を学習で知っている可能性があるので、予想には使わない。新しいモデルは、コードのレビューにだけ使う。
\end{itemize}

\section{対象の銘柄}

判断日に \SP{} に在籍し、実行 1 の予想と時価総額があり、過去 252 日の半分以上に株価がある銘柄（\code{blm/inputs.py} と同じ。平均約 488 銘柄）。どの戦略も、同じ判断日には同じ銘柄を使う。

\section{予想の当たり具合の指標}

\begin{longtable}{p{3.6cm}p{11.3cm}}
\toprule
\textbf{指標} & \textbf{定義} \\
\midrule
\endfirsthead
\bottomrule
\endfoot
IC（順位相関） & 判断日ごとに、予想の順位と翌月のリターンの順位のスピアマン相関。30 か月の平均と、$t$ 値（平均 $\div$ 標準偏差 $\times\sqrt{30}$） \\
相関（Pearson） & 判断日ごとの予想と翌月のリターンの相関の平均 \\
較正の傾き $b$ & 判断日ごとに平均を引いた予想と実際をまとめた回帰の傾き \\
MAE、RMSE & 予想（平均日次リターン）と翌月の平均日次リターンの差 \\
方向の正解率 & 判断日ごとの平均からの上下の向きが合った割合 \\
水準の偏り & 予想の平均 $-$ 実際の平均 \\
入力との相関 & 予想と、判断月の平均日次リターン（LLM が見た数字）の相関。なぞりの目安 \\
\end{longtable}

翌月のリターンは、判断日の終値から翌月の最後の取引日までの配当込みのリターン（ビュー \code{portfolio\_returns} と同じ作り方）を基本にする。資料の数値の再現（H01）では、資料と同じ「翌月の日次リターンの平均（15 日以上）」も使う。

\section{BL の挙動の指標}

採用率（銘柄ごとの $(\muBL-\Pi)/(q-\Pi)$ の中央値。相対ビューでは、全体の水準を除いた差で測る）、均衡からの乖離 $\|\muBL-\Pi\|$、$\tau\Sigma+\Omega$ の条件数、最適化の状態、能動的な比率のベータ $(w-\wmkt)\T\beta$。

\section{ポートフォリオの指標}

\begin{longtable}{p{3.6cm}p{11.3cm}}
\toprule
\textbf{指標} & \textbf{定義} \\
\midrule
\endfirsthead
\bottomrule
\endfoot
累積・年率リターン & 月次リターンの積から。年率は $(\prod(1+r_t))^{12/T}-1$ \\
年率ボラティリティ & $\mathrm{sd}(r_t)\sqrt{12}$ \\
シャープレシオ & 3 か月物国債（\code{\^{}IRX}）の利回りを月次に直して引いた超過リターンから \\
最大ドローダウン & 月次の資産の最高値からの下落率の最大 \\
TE、IR & 時価総額加重との月次の差の標準偏差 $\times\sqrt{12}$、差の平均 $\times12\,\div$ TE \\
$t$ 値 & 時価総額加重との月次の差の平均 $\div$ 標準誤差 \\
売買回転率 & 前月の比率を実際の値動きで動かした比率と、新しい比率の差の絶対値の和の半分 \\
コスト控除後 & 片道 0.1\%（第 2 版の論文と同じ）$\times$ 売買の量（回転率の 2 倍）を月次リターンから引く。比べる相手の時価総額加重にも同じ計算で引く \\
実効保有銘柄数 & $1/\sum_i w_i^2$ の平均 \\
ベンチマークからの乖離 & アクティブシェア $\frac12\sum_i|w_i-w_{\mathrm{mkt},i}|$ の平均 \\
\end{longtable}

\section{ポートフォリオの作り方（H03 以降で共通）}

結果を見る前に、次のとおり固定する。

\begin{itemize}
\item 共分散：過去 252 日の Ledoit--Wolf（現行のまま）。$\delta=2.5$、$\tau=0.025$（現行のまま）。
\item 最適化：$\max_w w\T\muBL-\frac{\delta}{2}w\T\Sigma w$、$\one\T w=1$、$w\ge0$ に、\textbf{事前の TE を年率 3\% 以内}（$(w-\wmkt)\T\Sigma(w-\wmkt)\le0.03^2/252$）を足す。3\% は『ビューとリスクのコツ』の計算と同じ値で、比べられるようにした。
\item 売買のコスト：片道 0.1\%。
\item 較正の傾き $b_t$：判断日より前に結果が分かっている月だけで計算し、$[0,\,0.99]$ に収める。3 か月分の結果がそろうまでは 0 とする。\textbf{$b_t=0$ の月は、ビューを使わず $\muBL=\Pi$ とする}（$\Omega=\tau\Sigma(1/b_t-1)$ は $b_t\to0$ で無限大になり、ビューの重みが 0 になる極限と同じ）。$\Omega\propto\Sigma$ の場合は、割り算を避けて同じ値になる $\muBL=\Pi+b_t\,\Sigma P\T(P\Sigma P\T)^{+}(q-P\Pi)$ で計算する。
\item 売買の時点：基本は「判断日の終値で売買」（資料と同じ）。ただし LLM の入力には判断日のリターンまで入っているので、厳密にはその終値では売買できない。そこで、\textbf{判断日の翌取引日の終値で売買する場合}も、主なポートフォリオについて感度分析として併記する。
\end{itemize}

\section{事前登録の範囲}

H02〜H04 は、既存の実行 1 の回答を使う再分析である。実行 1 の当たり具合の全体の値（IC $+0.05$ など）は、この設計の前に資料で知っていた。したがって H02〜H04 は「何も知らない状態での事前登録」ではなく、「全体の値を知ったうえで、新しい分析の指標と判定の基準を、その分析を行う前に固定したもの」である。H05〜H07 の LLM の回答は、まだ存在しない（本書のあとに初めて集める）。

\section{未来の情報が混ざらないことの確認}

\begin{itemize}
\item 財務の指標（\code{stock\_metrics}）は、判断日までに\textbf{提出された}四半期だけから作られている（テーブルの定義と資料の第 12 章。決算期末ではなく提出日で区切る）。150 日より古い決算は使わない。
\item 時価総額は、その日以前に報告された発行済み株式数と、その日の終値から作る。
\item 構成銘柄は、入れ替えの履歴から作った判断日時点の在籍期間で決める。データの取れない 28 銘柄の除外は、後の情報を使っている（資料の限界）。全戦略に共通なので比較は公平だが、成績の水準には偏りが残る。
\end{itemize}

\section{統計的な判断の決まり}

\begin{itemize}
\item 主な検定は、仮説ごとに 1 つの $p$ 値（両側）に決めておく（各章の「支持・棄却の条件」）。H02〜H07 の 6 つの主な検定には、Holm の方法で多重比較の補正をかける（全体の有意水準 5\%）。補正の前と後の両方を書く。
\item 30 か月の月次データで、IC の月ごとの標準偏差が約 0.19（資料の値）なら、$t$ 値 2 に届く IC の平均は約 0.07 である。\textbf{これより小さい IC は、本当にあっても検出できない。}
\item \textbf{有意でないことは、効果がないことの証明ではない。}「LLM の情報はない」の向きの主張には、意味のある最小の効果を \textbf{IC $=0.03$}（年率に直すと小さくない大きさで、資料の LLM の IC $0.05$ の 6 割）と決めておく。主な指標の 95\% 信頼区間の上限が 0.03 未満なら「効果はない（あっても小さい）」、上限が 0.03 以上で有意でないなら「結論を出すには証拠が不足している」と書く。
\item 順位を取るときの同じ値は、平均の順位にする。
\item わずかな数値の差や、1 か月だけの大きな成績は、一般化しない。
\end{itemize}

% ============================================================
\chapter{H01：現行の計算の再現と整合性の確認}\label{ch:h01}

\begin{plan}
仮説 ID・名称 & H01　現行の計算の再現と整合性の確認 \\
解決したい問題 & 研究の結論は、計算が正しいことに乗っている。資料の数値が今の DB とコードで再現できるか、式・単位・次元・収束に問題がないかを先に確かめる \\
資料との関係 & 総合報告書の問題 1〜6、ビューとリスクのコツ第 19 章の表を再現の対象にする \\
根拠 & 依頼文の優先順位 1。DB は 2015 年まで延長した版（\code{blm\_2015.sqlite}）に替わっており、資料を作った DB（\code{blm.sqlite}）と同じ結果になるかは確かめていない \\
変える要素 & なし（確認だけ）。テストのコードと、一般の $P$・行列の $\Omega$ を受け取る関数、TE つきの最適化、評価の仕組みを足す \\
固定する要素 & DB、実行 1 の回答、係数（$\tau=0.025$、$\delta=2.5$、252 日） \\
対照群 & 資料に書かれた数値 \\
データと期間 & 判断日 2024-03-28〜2026-08-31 の 30 か月 \\
評価指標 & \chapref{ch:math} の 7 項目 \\
支持・棄却の条件 & 次をすべて満たせば「整合している」とする。(1) 式~\eqref{eq:bl1} と式~\eqref{eq:bl2} の差が相対誤差 $10^{-8}$ 未満（ランダムな正定値行列と、実データ）。(2) ビューなしの比率と $\wmkt$ の差の最大が 0.1\% 未満。(3) Clarabel の状態がすべて最適。(4) 再現の値が資料の値と、相関・順位相関で $\pm0.01$、採用率で $\pm1$ ポイント以内。満たさないものは、原因を調べて記録する \\
実装と再現の手順 & \code{tests/}（標準ライブラリの \code{unittest}）と \code{experiments/H01\_verification/run.py}。\code{python -m unittest} と \code{python experiments/H01\_verification/run.py} で再現する \\
想定される失敗 & DB の作り直し（2015 年への延長）で、株価の調整後の値や銘柄の範囲が変わり、数値が少しずれる \\
API の費用 & 0 \\
\end{plan}

\begin{predbox}
式と単位は整合していて、資料の数値はほぼ再現できる。元の SLSQP は「成功」と報告しながら最適値の手前で止まっている。
\end{predbox}

% ============================================================
\chapter{H02：LLM の数値予想と統計モデル}\label{ch:h02}

\begin{hypbox}
LLM（実行 1）の翌月の予想が持つ情報は、LLM に渡した入力（判断月の平均日次リターンとその散らばり）の単純な関数で再現でき、それを統制すると、LLM に固有の予測力は残らない。また、LLM に渡していない数値（\code{stock\_metrics} の 12 の指標）から、評価期間より前のデータで学習した統計モデルの予想と比べて、LLM の予想は IC で上回らない。
\end{hypbox}

これは「LLM が統計モデルに対して追加の価値を生むか」を、最も素直な形で確かめるものである。帰無仮説を支持する方向の主張なので、結果が「支持されなかった」場合は、LLM が何かを足している証拠になる。

\begin{plan}
仮説 ID・名称 & H02　LLM の数値予想は、入力の単純な変換と統計モデルを超える情報を持つか \\
解決したい問題 & 相関 0.89 で入力をなぞる予想に、なぞり以外の情報があるかは確かめられていない \\
資料との関係 & 総合報告書の「統計モデルとの比較（未着手）」、ビューとリスクのコツ第 13 章の指標の IC（2015〜） \\
根拠 & 予想が入力の写しなら、写しを作る計算（回帰）で同じ予想が作れる。LLM に固有の情報は、写しとの差にしか残らない \\
変える要素 & 予想の作り方：(a) LLM（実行 1）、(b) LLM の写し（判断月の平均・散らばり・日数で LLM の予想を回帰した当てはめ値。ウォークフォワード）、(c) 短期反転（$-$直近 1 か月のリターン）、(d) モメンタム（12--1 か月）、(e) 統計モデル（12 の指標の断面 z スコアで、翌月のリターンの断面順位をリッジ回帰。2015-03〜の結果が確定した月で学習し、毎月学習し直す） \\
固定する要素 & 銘柄、期間、翌月のリターンの定義 \\
対照群 & (b)〜(e)。とくに (b) は「LLM を使わない同じ情報」の対照 \\
データと期間 & 評価は 30 か月。統計モデルの学習は 2015-03〜判断日の前 \\
評価指標 & IC と $t$ 値、相関、較正の傾き、MAE・RMSE、方向の正解率。写しの決定係数 \\
支持・棄却の条件 & \textbf{主な検定}：LLM の予想に固有の部分の IC。判断日ごとに、LLM の予想の順位を、判断月の平均日次リターンの順位と日次リターンの標準偏差の順位で回帰した残差（入力で説明できない部分）を作り、その残差と翌月のリターンの順位相関を取る。30 か月の平均の $t$ 値から両側の $p$ 値を出す。(i) $p<0.05$ で正なら棄却（LLM に固有の予測力がある）。(ii) 有意でなく、95\% 信頼区間の上限が 0.03 未満なら支持。(iii) 有意でなく、上限が 0.03 以上なら「証拠が不足」。（予想と入力の相関が 0.89 と高いので、両方を並べた回帰では標準誤差が大きくなる。そのため残差の IC を主な検定にし、並べた回帰は参考に示す。）\textbf{副次}：統計モデルと LLM の IC の差（月ごとの差の対応のある $t$ 検定） \\
実装と再現の手順 & \code{experiments/H02\_statistical\_baseline/run.py}。API は使わない \\
想定される失敗 & 統計モデルは、2015 年以降の指標の IC が小さい（資料）ので、統計モデルも当たらない可能性が高い。その場合「どちらも当たらない」という結果になり、比べる意味が弱くなる \\
API の費用 & 0 \\
\end{plan}

\begin{predbox}
写しの決定係数は 0.7 以上。残差の IC は有意にならない。ただし 30 か月では信頼区間の上限が 0.03 を下回らず、「証拠が不足」になる見込みが高い。統計モデルの IC も 0 と区別できない。
\end{predbox}

% ============================================================
\chapter{H03：ビューの形（$q$ と $P$）}\label{ch:h03}

\begin{hypbox}
LLM の絶対ビュー（$P=I$）には、全体の水準の楽観・悲観の癖が入っている。$\Omega$ が対角のとき、この癖は能動的な比率のベータの傾きとして比率に入る。予想を判断日ごとの時価総額加重の平均からの差（相対ビュー）にすれば、この傾きは消える。さらに、予想を順位の点数に直す（順位ビュー）と、極端な予想（上場直後の銘柄など）に比率が集中することを防げる。
\end{hypbox}

\begin{plan}
仮説 ID・名称 & H03　$q$ と $P$：絶対ビュー、相対ビュー、順位ビュー \\
解決したい問題 & LLM の予想は 76\% がプラスで、全体の水準の予想は翌月の実際と相関 $-0.43$（資料）。絶対ビューのままでは、当たらない水準の情報が比率に入る \\
資料との関係 & ビューとリスクのコツ第 9 章（全体の水準が比率に効くのは $\Omega$ が対角のとき）。第 19 章では $\Omega\propto\Sigma$ のもとで相対ビューと絶対ビューが同じ結果になることを確かめている。本仮説は、$\Omega$ が対角の場合と、順位ビューを新しく比べる \\
根拠 & 対角の $\Omega$ では、一律の上乗せ $c\one$ が $c\,\tau\Sigma(\tau\Sigma+\Omega)^{-1}\one$ になり、市場との共分散の大きい銘柄ほど大きく上がる \\
変える要素 & ビューの形 3 通り：(1) 絶対 $q=Q-r_f$、$P=I$。(2) 相対 $q_i=(Q_i-\wmkt\T Q)$、均衡の側も同じく平均を引く。(3) 順位：$z_i=\Phi^{-1}((\mathrm{rank}_i-0.375)/(N+0.25))$（同じ値は平均の順位）を、判断日ごとの予想の平均と標準偏差で日次の単位に戻した $q_i=\bar Q+\mathrm{sd}(Q)\,z_i-r_f$。順位だけを残し、外れ値の大きさを捨てる。3 通りとも同じ BL の式（同じ $\Omega$、同じ $b_t$）を通す。$\Omega$ は 2 通り：(a) 対角 $\tau\diag(\Sigma)(1/b_t-1)$、(b) $\Sigma$ に比例 $\tau\Sigma(1/b_t-1)$ \\
固定する要素 & 予想（実行 1）、$\Sigma$、$\delta$、$\tau$、TE 3\% の制約、コスト、銘柄、期間、$b_t$ の計算 \\
対照群 & 時価総額加重（ビューなし）。形の比較では (1) が対照 \\
データと期間 & 30 か月 \\
評価指標 & 能動的な比率のベータ、月ごとの予想の平均（楽観の度合い）と能動ベータの相関、最大の能動ウェイト、実効保有銘柄数、TE、IR、コスト控除後の超過リターン、回転率 \\
支持・棄却の条件 & \textbf{主な検定}：$\Omega$ が対角のとき、30 か月の「予想の時価総額加重の平均」と「能動的な比率のベータ」の相関の検定（相関 0 の $t$ 検定、両側）。絶対ビューで正かつ $p<0.05$ で、相対ビューの能動ベータの絶対値の平均が絶対ビューの半分以下なら支持。絶対ビューで有意でなければ「癖が比率に入る」こと自体が確かめられず、棄却。\textbf{副次}：順位ビューの最大の能動ウェイトが、絶対ビューより小さい（30 か月の平均） \\
実装と再現の手順 & \code{experiments/H03\_view\_shape/run.py}。API は使わない \\
想定される失敗 & TE 3\% の制約がベータの傾きを抑えてしまい、形の違いが見えにくい。その場合に備え、TE の制約なしでも同じ計算をして併記する \\
API の費用 & 0 \\
\end{plan}

\begin{predbox}
対角の $\Omega$ では、絶対ビューの能動ベータは楽観の月に大きくなる（相関が正）。相対ビューではそれが消える。$\Omega\propto\Sigma$ では形による違いはほぼない。成績の差は偶然と区別できない。
\end{predbox}

% ============================================================
\chapter{H04：$\Omega$ の設計と信頼度の較正}\label{ch:h04}

\begin{hypbox}
(1) 回答のばらつきは、銘柄の値動きの大きさを統制すると、予想の外れの大きさを説明しない。(2) $\Omega$ を判断日より前の当たり具合で較正し、$\Sigma$ に比例させ、推定の誤差に応じて縮めると、BL は設計どおりの採用率で「時価総額加重を少し動かす」形になる。(3) しかし、その比率の成績は、予想を銘柄の間でランダムに並べ替えたビュー（情報のない対照）と区別できない。
\end{hypbox}

\begin{plan}
仮説 ID・名称 & H04　$\Omega$：回答の分散、較正、縮小、置換による対照 \\
解決したい問題 & 回答の分散を $\Omega$ にすると、外れを 8 分の 1 に見積もり、予想を 91\% 採用する（資料）。較正の方法は資料で試されているが、(a) 推定の誤差に応じた縮小、(b) コスト、(c) 情報のない対照との比較、(d) 回答のばらつきが外れの順番の手がかりになるか、は確かめられていない \\
資料との関係 & ビューとリスクのコツ第 8、8b、19 章。そこでの 7 通りの比較を、コストつき・同じ TE の制約で再計算し、縮小と置換の対照を足す \\
根拠 & 傾きの推定値は少ない月では大きく揺れる（資料の図）。ランダムに並べ替えたビューで同じ成績が出るなら、成績はビューの情報ではなく、制約と銘柄の構成から来ている \\
変える要素 & $\Omega$ の 5 通り：(1) 回答の分散（対角）、(2) He--Litterman $\tau\diag(\Sigma)$、(3) 傾きで較正（対角）、(4) 傾きで較正（$\Sigma$ に比例）、(5) (4) の傾きを $\hat b\max\{0,1-\mathrm{se}^2/\hat b^2\}$ で縮めたもの（se は月ごとの傾きの標準誤差）。$\hat b\le0$ なら 0（ビューを使わない）。置換の対照：(5) と同じ $b_t$ で、予想を判断日ごとに銘柄の間で並べ替えたもの（種 0〜199 の 200 通り）。さらに主な比率について、判断日の翌取引日の終値で売買する場合も計算する \\
固定する要素 & ビューの形は相対ビュー、TE 3\%、コスト、$\Sigma$、$\delta$、$\tau$、銘柄、期間 \\
対照群 & 時価総額加重、置換のビュー \\
データと期間 & 30 か月 \\
評価指標 & 採用率、実効保有銘柄数、アクティブシェア、事前と実績の TE、IR、コスト控除後の超過リターン、回転率、ビューの成績の置換分布の中での位置（片側の経験的 $p$ 値） \\
支持・棄却の条件 & \textbf{主な検定}は (3)：(5) のコスト控除後の年率超過リターンの、置換の 200 通りの中での片側の経験的 $p$ 値（$(1+\#\{\text{置換}\ge\text{実際}\})/201$）。$p\ge0.05$ なら「ビューの情報は検出されない」（支持）、$p<0.05$ なら棄却。両側の Holm の補正には $2p$ を使う。(1) は、銘柄・月ごとに「回答の標準偏差」と「外れの絶対値」の順位相関を、過去 252 日のボラティリティで統制した偏順位相関で測り、30 か月の平均の $t<2$ なら支持。(2) は、(5) の実現の採用率が設定（$b_t$）から $\pm1$ ポイント以内、実効保有銘柄数が 30 以上、実績の TE が 6\% 以下なら支持 \\
実装と再現の手順 & \code{experiments/H04\_omega\_calibration/run.py}。較正の関数は \code{blm/calibration.py} \\
想定される失敗 & 計算の時間が長くなる（1 回の最適化 $\times$ 30 か月 $\times$ 約 205 通り）。時間が足りなければ置換の回数を減らし、その事実と $p$ 値の分解能を書く \\
API の費用 & 0 \\
\end{plan}

\begin{predbox}
(1) は支持（ばらつきは値動きの大きさで決まり、統制すると外れを説明しない）。(2) は支持。(3) も支持（LLM のビューの成績は置換と区別できない）。
\end{predbox}

% ============================================================
\chapter{H05：市場とセクターを、文章だけ・数字だけで聞く}\label{ch:h05}

\begin{hypbox}
実験 G で答えが市場の直近の値動きに沿ったのは、プロンプトに数字があったためである。運用会社の月次レポートの文章（要約と主なリスク）だけを渡すと、答えは市場の数字に沿わなくなる。ただし、その答えは翌月のセクターの成績を当てない。
\end{hypbox}

\begin{plan}
仮説 ID・名称 & H05　市場とセクター：文章だけ・数字だけで聞く（実験 G の拡張） \\
解決したい問題 & 「LLM を使う意味は文章を読む部分にある」という資料の結論の前提（文章から独自の判断が出るか）が、確かめられていない。資料が「次に行うべき実験」の 1 番目に挙げたもの \\
資料との関係 & 総合報告書の実験 G。G の 540 リクエスト（文章と数字の両方）の結果を、そのまま対照に使う \\
根拠 & G の点数は \SP{} の直近 1 か月のリターンと順位相関 $+0.77$、VIX と $-0.59$。文章だけにすれば、なぞる数字がない \\
変える要素 & 渡す情報：(G) 文章＋数字（既存）、(T) 文章だけ（日付、レポートの公表日、要約、主なリスク）、(N) 数字だけ（G の市場の数字とセクターの指標。レポートなし） \\
固定する要素 & 問い（リスク 7 種類、セクター 11 個）、答えの形（1〜5 の数字 1 文字）、システムプロンプト、gpt-4o-mini、温度 0 で確率を読む（G と同じ）、30 か月 \\
対照群 & (G)、(N) \\
データと期間 & 30 か月。文章は G のプロンプトに入っていたもの（レポートの読み取りの結果は今の DB にないため、G の結果のファイルから取り出す） \\
評価指標 & 点数（確率で重み付けた平均）の月ごとの散らばり、点数と \SP{} の直近 1 か月のリターン・VIX の順位相関、セクターの点数と翌月のセクター ETF の超過リターンの月ごとの順位相関（IC）、リスクの種類の符号の的中率 \\
支持・棄却の条件 & \textbf{主な検定}：(T) のセクターの IC の 30 か月の平均の $t$ 値（両側）。\textbf{なぞりの判定}：(T) の市場の点数と \SP{} の直近 1 か月のリターンの順位相関（30 か月）が、0 との差で有意でなく（$p\ge0.05$）、かつ (G) の値（$+0.77$）より有意に小さい（Fisher の $z$、$p<0.05$）なら「なぞりは消えた」。なぞりが消え、IC が正で $p<0.05$ なら「文章から当たる判断が出る」を支持。なぞりが消えて IC が有意でなければ「なぞりは消えるが、当たることは確かめられない」 \\
実装と再現の手順 & \code{experiments/H05\_text\_vs\_numbers/run.py}。回答は JSONL で保存 \\
想定される失敗 & 月次レポートは世界全体の見通しで、米国のセクターにほとんど触れない（資料：読み取れたのは最も多い情報技術で 39 本中 10 本）。(T) のセクターの答えが、ほぼ毎月同じになる可能性がある \\
API の費用 & 見込み：2 条件 $\times$ 540 リクエスト、入力約 700 トークン、出力 1 トークン（確率つき）。約 0.1〜0.2 ドル \\
\end{plan}

\begin{predbox}
(T) では市場の点数が直近のリターンに沿わなくなる。しかし月ごとの変化は小さく、セクターの IC は 0 と区別できない。
\end{predbox}

% ============================================================
\chapter{H06：決算発表文から LLM が取り出す情報}\label{ch:h06}

\begin{hypbox}
LLM に銘柄の数字（リターン、指標）を渡さず、各社の決算発表文（8-K の Item 2.02 に付く EX-99.1）だけを読ませて、業績の見通し・ガイダンスの変更・経営陣の語調を $-2$〜$+2$ の点数で取り出させる。この点数を、あらかじめ決めた等しい重みで足した合成点数は、翌月のリターンと正の順位相関を持ち、同じ時点で手に入る数値（売上と利益の前年比、直近 1 か月のリターン、12--1 か月のモメンタム）を統制しても、その相関が残る。
\end{hypbox}

これは依頼文のテーマ A の「LLM による情報抽出と、統計モデルによる数値化の分業」を、そのまま形にしたものである。LLM には数字を予想させず、文章からの判断だけをさせる。点数を期待リターンに直すのは、統計（Grinold の式とウォークフォワードの IC）である。

\begin{plan}
仮説 ID・名称 & H06　銘柄ごとの決算発表文から LLM が情報を取り出し、統計で数値化する \\
解決したい問題 & 月次レポートは全銘柄で同じ文章なので、銘柄の差を作れない（資料）。銘柄ごとの文章が必要 \\
資料との関係 & 総合報告書の次に行うこと 5（銘柄ごとの文章の情報源）、先行研究（Lopez-Lira と Tang、Anic ほか 2025。当初は Kim ほか 2024 も挙げたが、2025 年 2 月に取り下げられていたので外した。2026 年 10 月 19 日に訂正）。実験 D〜F は数字を渡したので、文章だけを渡すのは新しい \\
根拠 & 決算発表後の値動きは、発表の内容に沿って続く傾向（PEAD）が知られている。数字の驚きは数値でも測れるが、ガイダンスや語調は文章にしかない \\
変える要素 & ビューの元：(L) 決算発表文の合成点数。対照：(S) 同じ時点の数値の指標（売上と利益の前年比、モメンタム、直近 1 か月）、(N) 実行 1 の数値予想、(0) ビューなし \\
固定する要素 & gpt-4o-mini、温度 0、出力の形（JSON の構造化出力）、読む長さ（HTML のタグと表の罫線を取り除いた本文の先頭の約 15,000 文字）、点数の合成の重み（等しい重み、学習しない）、ポートフォリオの作り方（\chapref{ch:eval}） \\
対照群 & (S)、(N)、(0) \\
データと期間 & 決算発表文：SEC EDGAR から、提出日が 2023-11-01〜2026-08-30 のもの。判断日ごとに、提出日が判断日より前で 120 日以内の最新の 1 本を使う。gpt-4o-mini の知識の期限（2023 年 10 月）より後の文章だけである \\
評価指標 & IC と $t$ 値（全銘柄、発表から 35 日以内の銘柄）、Fama--MacBeth の係数、点数と直近リターンの相関（なぞりの確認）、ポートフォリオの指標 \\
支持・棄却の条件 & \textbf{主な検定}：判断日ごとに、翌月のリターンの順位を、合成点数の順位と、売上の前年比・利益の前年比・直近 1 か月のリターン・12--1 か月のモメンタムの順位で回帰し（Fama--MacBeth）、合成点数の係数の平均の $t$ 値（両側）。正で $p<0.05$ なら支持、有意でなければ「支持されなかった」（信頼区間も示す）。副次：合成点数の IC の $t$、発表から 35 日以内の銘柄だけでの IC \\
実装と再現の手順 & \code{experiments/H06\_earnings\_text/}。取得（\code{fetch.py}）、読み取り（\code{score.py}）、評価（\code{evaluate.py}）の 3 段。取得した文章は \code{data/} の下（Git の外）に置き、点数と条件は実験のフォルダに保存する。SEC へのリクエストは毎秒 10 回以下に抑え、連絡先は実行時の環境変数で渡す \\
想定される失敗 & 発表文が EX-99.1 でない会社がある（ファイル名で見つける規則が合わない）。文章が長く、先頭だけでは見通しの部分が切れる。点数がほとんどの会社で同じ（例：語調は毎回「やや前向き」）になり、銘柄の差が出ない。30 か月では小さな IC を検出できない \\
API の費用 & 見込み：約 5,500 本 $\times$ 入力約 4,500 トークン $\approx$ 2,500 万トークンで約 3.7 ドル、出力約 0.5 ドル。合計約 4〜5 ドル。最初に 30 本で動作を確かめる \\
\end{plan}

\begin{predbox}
点数は直近のリターンとはあまり相関しない（なぞりは起きにくい）。IC は小さな正（0.01〜0.03）で、30 か月では $t<2$ にとどまる可能性が高い。発表から日が浅い銘柄だけでは IC が大きくなる。
\end{predbox}

% ============================================================
\chapter{H07：フィードバックの LLM}\label{ch:h07}

\begin{hypbox}
LLM に自分の過去の予想の外れを渡すと、予想の水準の偏り（楽観）は小さくなる。しかし、銘柄の順位の当たり具合（IC）は上がらない。銘柄ごとの前月の外れを渡すと、外れを追いかけて予想を大きく動かす（過剰な反応）。予想する役とフィードバックを受けて直す役を分けると、予想は小さく縮むが、順位は変わらない。
\end{hypbox}

既存のコードには、フィードバックの LLM がない。そこで、何を渡すと何が改善するはずかを先に決め、小さな規模で比べる。

\begin{center}
\small
\begin{tabular}{p{2.6cm}p{6.2cm}p{5.4cm}}
\toprule
条件 & 渡すもの & 改善するはずのもの（理由） \\
\midrule
F0 フィードバックなし & 実行 1 と同じプロンプト（実行 1 の回答をそのまま使う）。雑音の目安として、同じプロンプトでもう 1 回聞く（F0$'$） & --- \\
F1 外れのみ & F0 ＋「先月のあなたの予想は $x$、実際は $y$」（その銘柄） & 銘柄ごとの癖が直るはず。ただし 1 か月の外れは雑音が大きい \\
F2 較正の情報 & F0 ＋ 判断日までの全銘柄の成績の要約（予想の平均と実際の平均、順位相関、入力との相関） & 水準の偏りと、なぞりの自覚 \\
F3 役の分離 & 別の呼び出しで「見直す役」に、F0 の予想の平均、入力、F2 と同じ要約を渡し、直した予想を答えさせる & 予想する役の癖を、別の役が割り引く \\
\bottomrule
\end{tabular}
\end{center}

「BL を適用した後の比率やリスクを渡す」条件は、この規模では行わない。理由は、相対ビューと $\Omega\propto\Sigma$ のもとでは水準の偏りが比率に効かない（H03、H04 で確かめる）ので、比率の情報で予想の何が良くなるべきかを、事前に決められないためである。結果の報告書で、この判断を見直す。

\begin{plan}
仮説 ID・名称 & H07　フィードバックの LLM：何を渡すと何が変わるか \\
解決したい問題 & フィードバックの LLM は未実装。フィードバックを足すことで、過学習、自己強化、不要な予想の変更が起きないかも見る \\
資料との関係 & 総合報告書の段階 2（未着手）。実験 F（会話にすると方針が固まり、なぞりが強まった） \\
根拠 & 予想の外れの情報は判断日の時点で手に入る（判断日 $d_{i-1}$ の予想の結果は $d_i$ で確定する）ので、未来の情報を使わずに渡せる \\
変える要素 & F0〜F3 の 4 条件（上の表） \\
固定する要素 & gpt-4o-mini、温度 1.0、1 リクエストで 30 回答、F0 と同じ入力（判断月の日次リターンの列と会社情報）、銘柄、期間 \\
対照群 & F0（実行 1）。雑音の目安として F0$'$ \\
データと期間 & 最初の判断日（2024-03-28）に実行 1 の予想がある銘柄から、種 0 で 100 銘柄を選ぶ（その後の在籍を条件にすると未来の情報を使うので、条件にしない）。指数から外れた銘柄は、その月から除き、補わない。30 か月（フィードバックは 2 か月目から） \\
評価指標 & IC、較正の傾き、水準の偏り、入力との相関、F0 からの予想の変化の大きさ（F0$'$ と F0 の差に対する倍率）、回答の標準偏差 \\
支持・棄却の条件 & \textbf{主な検定}：F2 と F0 の月ごとの IC の差の対応のある $t$ 検定（両側）。正で $p<0.05$ なら棄却（フィードバックで IC が上がる）。有意でなく、差の 95\% 信頼区間の上限が 0.03 未満なら支持、上限が 0.03 以上なら「証拠が不足」。副次：水準の偏りの絶対値が F0 の半分以下になるか（F2、F3）、F1 の予想の変化が雑音の目安の 2 倍を超えるか \\
実装と再現の手順 & \code{experiments/H07\_feedback\_llm/run.py}。判断日の順に 1 か月ずつ実行する。全回答、プロンプト、トークン数を JSONL で保存 \\
想定される失敗 & 100 銘柄 $\times$ 29 か月では IC の推定の誤差が大きい（月ごとの IC の標準偏差は、銘柄数が少ないぶん大きくなる） \\
API の費用 & 見込み：4 条件（F0$'$、F1、F2、F3）$\times$ 約 2,900 リクエスト、入力約 350 トークン、出力 30 回答 $\times$ 約 8 トークン。約 2〜3 ドル \\
\end{plan}

\begin{predbox}
F2 と F3 は水準の偏りを小さくする。IC は F0 と区別できない。F1 は予想を雑音の目安より大きく動かすが、IC は上がらない。
\end{predbox}

% ============================================================
\chapter{リスクと限界}\label{ch:risk}

\begin{longtable}{p{4.0cm}p{10.9cm}}
\toprule
\textbf{リスク・限界} & \textbf{内容と対処} \\
\midrule
\endfirsthead
\bottomrule
\endfoot
評価の期間が短い & 30 か月では、IC が 0.07 未満の効果は検出できない。「支持されなかった」と「効果がない」を区別して書く \\
多重比較 & 仮説が 7 つ、各仮説に複数の変種がある。主な検定を仮説ごとに 1 つに決め、Holm の補正をかける。副次の結果は参考として扱う \\
相場の局面の偏り & 評価期間は上昇相場が中心で、下落相場の挙動は分からない \\
生存者バイアス & データの取れない 28 銘柄の除外（資料）は全戦略に共通。成績の水準には偏りが残る \\
LLM の乱数 & 温度 1.0 の実験は、同じ条件でも回答が変わる。全回答と条件を保存し、「保存した回答から同じ結果を計算し直せる」ことを再現性とする（資料の決まりと同じ） \\
LLM の知識の期限 & gpt-4o-mini（2023 年 10 月まで）だけを予想に使う。決算発表文には社名が入るが、評価期間の出来事はモデルが知らない \\
決算発表文の取得 & SEC の利用の決まり（毎秒 10 回以下、連絡先の明記）を守る。取得した文章は Git に入れない \\
コードのレビュー & 新しいモデル（gpt-5.x など）にコードを読ませるのは、レビューの補助だけにする。その指摘も、自分で確かめてから使う \\
費用 & 見込みの合計は 10 ドル前後（総合報告書の予算の上限 5 万円の範囲）。各実験の前に 30 件程度で動作を確かめる \\
\end{longtable}

\section*{結果の報告で守ること}

\begin{itemize}
\item 仮説ごとの結論は「支持された」「支持されなかった」「結論を出すには証拠が不足している」「実験未実施」のどれかで書く。
\item 成績が改善しなかった場合も、そのまま記録する。対照群との差を数字で書く。
\item 実行できなかった部分は、実行した部分と分けて書き、推測で埋めない。
\end{itemize}

\end{document}
